% Point to the main file
\documentclass[../../Main.tex]{subfiles}

% ========== Start the document ==========
\begin{document}

\section{Section 6.4 - Binomial Identities}

\begin{align*}
    (a + b)^0 &= 1 \\
    (a + b)^1 &= 1a + 1b \\
    (a + b)^2 &= 1a^2 + 2a^1b^1 + 1b^2 \\
    (a + b)^3 &= 1a^3 + 3a^2b^1 + 3a^1b^2 + 1b^3 \\
    (a + b)^4 &= 1a^4 + 4a^3b^1 + 6a^2b^2 + 4a^1b^3 + 1b^4 \\
    (a + b)^5 &= 1a^5 + 5a^4b^1 + 10a^3b^2 + 10a^2b^3 + 5a^1b^4 + 1b^5 \\
    &\vdots \\
    (a + b)^{10} &= 1a^{10} + \underbrace{10}_{10 \choose 9}a^9b^1 + \underbrace{45}_{10 \choose 8}a^8b^2 + \underbrace{120}_{10 \choose 7}a^7b^3 + \underbrace{210}_{10 \choose 6} a^6b^4 + \underbrace{252}_{10 \choose 5}a^5b^5 + \underbrace{210}_{10 \choose 4}a^4b^6 + 120a^3b^7 + 45a^2b^8 + 10 a^1b^9 + 1b^10 \\
\end{align*}

For all of these, it's 10 choose either the power of $a$ or the power of $b$. This is because the power is essentially the number of arrangements of $aaabbbbbbb$ for example

This brings us to our next theorem

\subsection{Binomial Theorem}

\begin{theorem}[Binomial Theorem]
    
    $$\text{For } n = 0, 1, 2, \ldots, (a + b)^n = \sum_{k = 0}^{n} {n \choose k} a^{n - k}b^k$$

    I would put the proof here but it's long and I'm lazy

\end{theorem}

\begin{example}
    \begin{enumerate}
        \item 
        {
            A pizza parlor has 10 toppings. How many pizzas are possible if any number of toppings can be chosen but none can be repeated?

            \begin{align*}
                \underbrace{\text{\# 0 toppings}}_{10 \choose 0} &+
                \underbrace{\text{\# 1 topping}}_{10 \choose 1} +
                \ldots +
                \underbrace{\text{\# 0 toppings}}_{10 \choose 10} \\
                &= \sum_{k = 0}^{10} {10 \choose k} \underbrace{1^{10 - k} \cdot 1^k}_{\text{to use B.T}} \\
                &= (1 + 1)^{10} \\
                &= 2^{10} \\
                &= \boxed{1024} \\
            \end{align*}
        }

        \item
        {
            How many bit strings are there of length 10?

            $$\underline{2} \cdot \underline{2} \cdot \underline{2} \cdot \underline{2} \cdot \underline{2} \cdot \underline{2} \cdot \underline{2} \cdot \underline{2} \cdot \underline{2} \cdot \underline{2} = 2^{10} = \boxed{1024}$$

            or

            Place $k$ $0$'s in the bit string.

            $${10 \choose 0} + {10 \choose 1} + {10 \choose 2} + \ldots + {10 \choose 10} = (1 + 1)^{10} = 2^{10} = 1024$$
        }

        \item
        {
            How many subsets does a set with 10 elements have?

            uhhhh same thing
        }
    \end{enumerate}
\end{example}

\subsection{Sum of Combinations}

\begin{theorem}[Sum of Combinations]
    $$\text{For } n = 0, 1, 2, \ldots, \sum_{k = 0}^{n}{n \choose k} = 2^n$$
\end{theorem}

\begin{example}
    uhhhhh just use the formula
\end{example}

% ========== End the document ==========
\end{document}